\part{Reading the code}

\chapter{The spine: one tool call, end to end}

The fastest way into \nodwyn{} is to follow a single tool call from the
keyboard to the disk and back. Eleven files, in order. The excerpts below
are the real code, lightly abridged; line numbers refer to \code{v0.1}.

\section{The front-end sends a line}

\file{internal/frontend/plain/chat.go}. The plain front-end reads a line,
hands it to the session, and then does nothing but render events until the
session says it is idle. It never touches the loop, the guard or the
interpreter\,---\,it is a log subscriber with a reply channel.

\begin{go}
if err := c.Session.Send(line); err != nil {
    fmt.Fprintln(c.Out, "error:", err)
    continue
}
c.follow(ctx)

// follow renders events until the turn ends.
func (c *Chat) follow(ctx context.Context) {
    for {
        select {
        case <-c.Interrupt:
            fmt.Fprintln(c.Out, "\n[aborting]")
            c.Session.Abort()
        case ev, ok := <-c.Session.Events():
            if !ok {
                return
            }
            c.show(ev)
            if ev.Kind == event.Status &&
                ev.Payload.(event.StatusInfo).State == "idle" {
                return
            }
        }
    }
}
\end{go}

\section{The session starts a turn}

\file{internal/core/session.go}. \code{Send} refuses if a turn is running,
otherwise starts one in a goroutine with a cancellable context. When the
turn ends\,---\,however it ends\,---\,the session emits \code{status idle}. Note
\code{turnTrust} being reset: a ``trust writes this turn'' answer cannot
outlive the turn.

\begin{go}
func (s *Session) Send(text string) error {
    s.mu.Lock()
    if s.busy {
        s.mu.Unlock()
        return errors.New("core: a turn is already running")
    }
    ctx, cancel := context.WithCancel(context.Background())
    s.busy, s.turnCancel = true, cancel
    s.mu.Unlock()

    go func() {
        defer func() {
            cancel()
            s.mu.Lock()
            s.busy, s.turnCancel, s.turnTrust = false, nil, false
            s.mu.Unlock()
            s.Emit(event.Event{Kind: event.Status,
                Payload: event.StatusInfo{State: "idle"}})
        }()
        if err := s.loop.Turn(ctx, text); err != nil {
            s.Emit(event.Event{Kind: event.Error,
                Payload: event.ErrorInfo{Err: err.Error()}})
        }
    }()
    return nil
}
\end{go}

\code{Record} is the one function that writes the log, keeps the event for
the window and forwards it to the front-end. It is what the loop, the guard
and the session itself all call:

\begin{go}
func (s *Session) Record(kind log.Kind, payload any) (log.Event, error) {
    ev, err := s.logw.Append(kind, payload)
    if err != nil {
        return ev, err
    }
    s.evmu.Lock()
    s.events = append(s.events, ev)
    s.evmu.Unlock()
    s.Emit(event.FromLog(ev, payload))
    return ev, nil
}
\end{go}

\section{The loop asks the model}

\file{internal/loop/loop.go}. The loop itself is Scheme: \code{Turn} applies
\code{agent-turn}, the default from \file{internal/loop/strategy.scm} or a
project's redefinition, with harness origin, then answers any call the
strategy left unanswered and closes the turn's budget.

\begin{go}
func (l *Loop) Turn(ctx context.Context, request string) error {
    strategy, _ := l.interp.Global().Lookup("agent-turn")
    l.turn = &turnState{budget: budget.From(ctx)}
    defer func() { l.turn = nil }()
    _, err := l.interp.Apply(scheme.WithOrigin(ctx, scheme.OriginHarness),
        strategy, []scheme.Value{scheme.Str(request)})
    if skipErr := l.skipPending("the turn ended before this call ran"); skipErr != nil && err == nil {
        err = skipErr
    }
    if err != nil {
        return surfaced(ctx, err)
    }
    return l.turn.budget.Finish(ctx)
}
\end{go}

The default strategy asks, runs what the model asked for, and asks again
until a reply has no tool calls:

\begin{scm}
(define (agent-turn request)
  (let loop ((reply (model-turn)))
    (let ((calls (assoc-ref reply 'tool-calls '())))
      (unless (null? calls)
        (run-tools calls)
        (maybe-compact request reply)
        (loop (model-turn))))))
\end{scm}

\code{model-turn} (\file{internal/loop/model.go}) is where the model is
asked. Every request rebuilds the window from the events so far (there is no
``messages'' field kept between requests), records an \code{llm.req}
describing the window, streams the response, and records the full
\code{llm.resp}. The budget, the tool-call limit and the progress guard sit
in the primitives, where no strategy can skip them.

\section{The window is a fold over the log}

\file{internal/window/window.go}. \code{Build} turns \code{user.msg},
\code{llm.resp} and \code{tool.result} events into messages and applies a
policy shaped for a prefix cache on one shared model slot: the system
prompt carries the \emph{current task}, so the ask is never trimmed;
elision is \emph{sticky}\,---\,every earlier request's elided set, read back
from the log's \code{llm.req} events, stays elided; and \emph{hysteresis}
trims above a high-water mark down to a low-water mark, oldest first,
tool results before assistant text, user messages never, then holds until
the next crossing. Between trims every request is a byte-identical prefix
plus a delta.

\begin{go}
if total() > pol.HighWater {
    trimmed = true
    protect := len(entries) - Keep
    // Pass 1: oldest tool results. Pass 2: oldest assistant turns.
    for _, kind := range []string{"tool", "assistant"} {
        for i := 0; i < protect && total() > pol.LowWater; i++ {
            if entries[i].kind == kind && !isElided(i) {
                elide(i)
            }
        }
    }
}
\end{go}

\section{The provider streams}

\file{internal/llm/ollama.go}. The native \code{/api/chat} protocol sends
newline-delimited JSON chunks. Content and thinking are accumulated and also
forwarded as deltas; tool calls arrive as JSON objects and are kept as such;
the final chunk carries usage. Nothing in the loop knows this is Ollama.

\begin{go}
for sc.Scan() {
    var ch wireChunk
    if err := json.Unmarshal(sc.Bytes(), &ch); err != nil {
        return nil, fmt.Errorf("ollama: bad chunk: %w", err)
    }
    if ch.Message.Content != "" || ch.Message.Thinking != "" {
        content.WriteString(ch.Message.Content)
        thinking.WriteString(ch.Message.Thinking)
        onDelta(Delta{Text: ch.Message.Content, Thinking: ch.Message.Thinking})
    }
    for _, wt := range ch.Message.ToolCalls {
        out.Message.ToolCalls = append(out.Message.ToolCalls,
            ToolCall{ID: wt.ID, Name: wt.Function.Name, Args: wt.Function.Arguments})
    }
    if ch.Done {
        out.DoneReason = ch.DoneReason
        out.Usage = Usage{PromptTokens: ch.PromptEvalCount, /* ... */}
        break
    }
}
\end{go}

\section{A tool call becomes a Scheme application}

\file{internal/loop/loop.go}, \code{apply}. This is \adr{3} in four lines: look
the primitive up, order the model's named JSON arguments into the
primitive's positional parameters (unknown or missing names are errors the
model is told about), tag the context with \emph{model origin}, and apply.
The model's call now takes exactly the path a call from the REPL would.

\begin{go}
func (l *Loop) apply(ctx context.Context, tc llm.ToolCall) (string, error) {
    p, ok := l.interp.Lookup(tc.Name)
    if !ok || !p.Tool() {
        return "", fmt.Errorf("unknown tool %q; available: %s", tc.Name, l.toolNames())
    }
    args, err := p.ArgsFromJSON(tc.Args)
    if err != nil {
        return "", err
    }
    v, err := l.interp.Apply(scheme.WithOrigin(ctx, scheme.OriginModel),
        scheme.Value{Kind: scheme.KPrim, Prim: p}, args)
    if err != nil {
        return "", err
    }
    return render(v), nil
}
\end{go}

\section{The single exit}

\file{internal/scheme/prim.go}. Whatever the caller\,---\,the loop, \code{map},
\code{apply}, a registered tool's body, the REPL\,---\,every primitive
application reaches this function and no other. Arity first, then the gate,
then the handler. A lint asserts that \code{p.Handler} is invoked nowhere
else.

\begin{go}
// applyPrim is the single Scheme->Go exit. Every primitive application --
// from the model, the prelude, a .scm file or the REPL -- passes through
// here: arity check, guard, handler. Nothing else calls p.Handler.
func (in *Interp) applyPrim(ctx context.Context, p *Prim, args []Value) (Value, error) {
    if err := ctx.Err(); err != nil {
        return Value{}, wrap(p.Name, err)
    }
    if err := checkArity(p, args); err != nil {
        return Value{}, err
    }
    if in.gate != nil {
        if err := in.gate.Before(ctx, p, args); err != nil {
            return Value{}, wrap(p.Name, err)
        }
    }
    v, err := p.Handler(&Call{Ctx: ctx, In: in, Name: p.Name, Args: args,
        Origin: OriginFrom(ctx)})
    if err != nil {
        return Value{}, wrap(p.Name, err)
    }
    return v, nil
}
\end{go}

\section{The guard decides}

\file{internal/guard/guard.go}. Checkers run in order; a deny short-circuits;
the strictest verdict wins; a standing approval converts a confirm into an
allow (and carries the diff into the log, since nobody saw it). A confirm
goes to the confirmer, which the core implements as an event plus a wait.

\begin{go}
for _, c := range g.checkers {
    d := c.Check(ctx, p, args)
    if d.Verdict == Allow {
        if d.Standing && standing == "" {
            standing = d.Reason
        }
        continue
    }
    reasons = append(reasons, d.Reason)
    if d.Detail != "" {
        details = append(details, d.Detail)
    }
    if d.Verdict > final.Verdict {
        final, checker = d, c.Name()
    }
    if d.Verdict == Deny {
        break // nothing can relax a deny
    }
}
if final.Verdict == Confirm && standing != "" {
    final = Decision{Verdict: Allow, Standing: true, Reason: standing}
    reasons, checker = []string{standing}, "standing"
}
\end{go}

\section{A checker reads metadata, not names}

\file{internal/guard/checkers.go}. \code{RepoPin} does not know what
\code{fs-read} is. It asks the primitive's declaration which parameters
carry a path, resolves each against the root\,---\,following symlinks on the
longest existing prefix, so a link out of the tree does not escape\,---\,and
denies anything that lands outside.

\begin{go}
func (r RepoPin) Resolve(raw string) string {
    abs := raw
    if !filepath.IsAbs(abs) {
        abs = filepath.Join(r.Root, abs)
    }
    abs = filepath.Clean(abs)
    resolved := resolveExisting(abs)
    root := resolveExisting(r.Root)
    if resolved == root || strings.HasPrefix(resolved, root+string(filepath.Separator)) {
        return abs
    }
    return ""
}
\end{go}

\section{The primitive runs}

\file{internal/prim/fs.go}. \code{fs-edit} replaces exactly one occurrence.
Its \code{Preview} renders the diff the confirm prompt shows; its handler
re-resolves the path (defense in depth: the guard already denied an escape)
and refuses ambiguous or missing matches with messages written for the
model to act on.

\begin{go}
func applyEdit(content, old, replacement string) (string, error) {
    if old == "" {
        return "", errors.New("old text must not be empty")
    }
    switch n := strings.Count(content, old); n {
    case 0:
        return "", errors.New("old text not found; it must match exactly, including whitespace")
    case 1:
        return strings.Replace(content, old, replacement, 1), nil
    default:
        return "", fmt.Errorf("old text occurs %d times; include more context so it is unique", n)
    }
}
\end{go}

\section{The result is recorded, and shown}

Back in \file{loop.go}, \code{runTool} records \code{tool.call} before and
\code{tool.result} after. Only context cancellation is an error here;
everything else\,---\,a denial, a missing file, a bad argument\,---\,is a
\emph{result} the model reads:

\begin{go}
res := log.ToolResult{ID: tc.ID, Name: tc.Name}
value, err := l.apply(ctx, tc)
if err != nil {
    if ctx.Err() != nil {
        return ctx.Err()
    }
    res.OK, res.Error = false, err.Error()
} else {
    res.OK, res.Value = true, value
}
_, err = l.host.Record(log.KindToolResult, res)
\end{go}

And \file{internal/log/log.go} writes it: marshal, filter (the future
redactor), one line, fsync.

\begin{go}
line, err := json.Marshal(ev)
line = append(line, '\n')
if _, err := w.w.Write(line); err != nil {
    return Event{}, fmt.Errorf("log: write: %w", err)
}
if w.syncer != nil {
    if err := w.syncer.Sync(); err != nil {
        return Event{}, fmt.Errorf("log: fsync: %w", err)
    }
}
w.seq = ev.Seq
\end{go}

The front-end, meanwhile, has been receiving each of these as an event and
printing a tool card, a guard flag, or a dimmed diff. That is the whole
trip.

\chapter{The log as truth}

A session log lives under \file{\textasciitilde/.local/share/nodwyn/sessions/}, one \file{<id>.jsonl} per session. The
kinds are fixed in \file{internal/log/log.go}; a reader rejects any other,
and rejects a gap in \code{seq}. Three lines from a real session:

\begin{txt}
{"seq":2,"ts":"2026-09-21T03:14:52Z","session":"20260921T031452Z-da2e39d8",
 "kind":"user.msg","payload":{"text":"Read internal/log/log.go and explain ..."}}
{"seq":5,"ts":"...","kind":"tool.call","payload":{"id":"call_x","name":"fs-read",
 "args":{"path":"internal/log/log.go"},"origin":"model"}}
{"seq":6,"ts":"...","kind":"guard.decision","payload":{"prim":"fs-read",
 "origin":"model","decision":"allow","args":"(\"internal/log/log.go\")"}}
\end{txt}

What is \emph{not} in the log is as telling as what is: \code{llm.req} records
how many messages and an estimate, not the messages\,---\,they are
reconstructible. Model thinking is stored on \code{llm.resp} and never sent
back. Pure builtins (\code{car}, \code{+}) are allowed silently; anything that
reaches the host, and every non-allow, is recorded.

Two properties come for free once the log is the truth. A crash between
write and fsync leaves a partial last line; \code{Open} repairs it and says
so. And ``resume'' (not built yet) is \code{Build} over an old file.

\chapter{Conventions}

\paragraph{Metadata is the contract.} A primitive declares \code{Doc},
\code{Params} (name, type, doc, \emph{role}), \code{Variadic}, \code{Egress},
\code{Irreversible}, \code{Host}, \code{Hidden} and an optional \code{Preview}.
The model's tool schema and the guard's judgement both come from this one
declaration and cannot drift. Roles\,---\,\code{path}, \code{argv},
\code{payload}, later \code{dest} and \code{credential}\,---\,are what checkers
read. A checker that switches on a primitive's name is a bug.

\paragraph{Origins.} A context carries who initiated the evaluation:
\code{user} (REPL, \code{:eval}, \file{init.scm}) or \code{model} (a tool
call). Every guard record has it; L2 will use it to restrict tag stripping.

\paragraph{Errors.} \code{*scheme.Error} is the Scheme-level error, with an
optional \code{Cause} for host errors (a \code{*Denied}, a cancelled
context). Runtime failures read \code{who: what}. A tool's failure is a
result for the model; only cancellation is an error for the loop.

\paragraph{Two binaries.} \file{bin/nodwyn} is what you are editing;
\code{nodwyn-stable} is what edits it. \code{make promote} is the only bridge.

\paragraph{One ingress, one exit.} There is one function that writes the
log, one that appends to the model's context (the window), one that calls a
primitive's handler. Lints and tests hold each. When a phase needs a second
one, the answer is an ADR, not a shortcut.
